\chapter{Python and the Testbench}
\label{ch:python}

\chapabstract{Your golden model is easier to write in Python than in \SV{} or
\VHDL{}, and you already know how to write it. This chapter answers the
question that follows: how do you connect it to a simulation? It gives a
taxonomy of four coupling styles -- offline vectors, generated source,
in-process calls, and out-of-process co-simulation -- and works through each
one with code that has actually been run against the CORDIC rotator. At the
end you will be able to pick the right style for a homework exercise, a
thesis accelerator or an industrial block, and to build the flow that goes
with it.}

\section{Why Python belongs in a verification flow}
\label{sec:python:why}

A CORDIC reference model is three lines of \code{numpy}. A fixed point
reference model of the same rotator, bit for bit, is about thirty lines of
Python and about a hundred lines of \SV{}. That asymmetry is the whole
argument, and it is worth being precise about where it comes from rather than
repeating it as a slogan.

Python is good at five things a verification engineer needs:

\begin{itemize}
\item \textbf{Golden models.} Arbitrary-precision integers, real
      arithmetic, \code{numpy} arrays and \code{scipy} transforms are all one
      import away. Nothing in a Python model has a width, so the model never
      accidentally truncates. You add the truncation deliberately, which is
      exactly what a fixed point reference model is: a floating point
      algorithm plus a set of explicit rounding and saturation decisions.
\item \textbf{Stimulus generation and corner-case enumeration.} Generating
      every angle that lands exactly on a quadrant boundary is a list
      comprehension. Enumerating the \(2^{16}\) angles of a 16-bit input and
      filtering the interesting ones takes one line and a second of CPU time.
\item \textbf{Post-processing.} A simulation produces numbers. Deciding
      whether those numbers are acceptable often means computing an error
      spectrum, a histogram or a plot -- work that no HDL is designed for.
\item \textbf{Code generation.} Register maps, ROM contents, decode tables,
      testbench skeletons, constant packages. Anything you would otherwise
      compute by hand and paste into a file.
\item \textbf{Glue.} Makefiles, CI jobs, log parsing, result aggregation
      across a regression.
\end{itemize}

\index{golden model}\index{reference model}
None of this is new. Verification teams have used Perl, Tcl and C for the same
purposes for thirty years. What changed is that the numerical libraries moved
to Python, and with them the people who write DSP and machine-learning
hardware.

\subsection{The honest counter-argument}

The in-language model is sometimes the better choice, and a chapter that
pretends otherwise would be selling something. \SV{} has \kw{real}
arithmetic, the \code{\$atan}/\code{\$sin}/\code{\$exp} family of real-valued
system functions, and, since IEEE 1800-2005, \code{shortreal} as well.
\VHDL{} has \code{ieee.math\_real}, which is older and arguably better
specified. Both are perfectly capable of expressing a reference model.

\index{math\_real}\index{\$atan}
Here is the arctangent table of the CORDIC computed at elaboration time in
each language. The left pane is the real \file{cordic\_pkg.vhd}; the right
pane is the real \file{cordic\_pkg.sv}.

\begin{rosetta}
\begin{rleft}
use ieee.math_real.all;

type atan_rom_t is
  array (0 to N-1) of iangle_t;

function build_atan_rom
  return atan_rom_t is
  variable r : atan_rom_t;
begin
  for i in 0 to N-1 loop
    r(i) := to_signed(integer(round(
      arctan(2.0 ** (-i))
        * 2.0 ** IQA)), IA);
  end loop;
  return r;
end function;

constant ATAN_ROM : atan_rom_t
  := build_atan_rom;
\end{rleft}
\begin{rright}
typedef iangle_t atan_rom_t [N];

function automatic atan_rom_t
    build_atan_rom();
  atan_rom_t rom;
  for (int i = 0; i < int'(N); i++)
    rom[i] = iangle_t'($rtoi(
      $atan(2.0 ** (-i))
        * (2.0 ** IQA) + 0.5));
  return rom;
endfunction

localparam atan_rom_t ATAN_ROM =
  build_atan_rom();
\end{rright}
\end{rosetta}

The two are the same idea: a pure function, called once by the tool, whose
result becomes a constant. \VHDL{} has \code{round}, so the intent is
explicit; \SV{} has only \code{\$rtoi}, which truncates towards zero, so the
\code{+ 0.5} is doing the rounding by hand. That is the kind of detail that
decides whether your model and your \RTL{} agree in the last bit.

When the model is this small, keeping it in the \RTL{} language is better:
there is one source file, no build step, no generated artefact to go stale,
and the constants are visible in the waveform viewer under their own names.
Reach for Python when at least one of the following is true.

\begin{ruleofthumb}
Write the model in the \RTL{} language if it fits on one screen and uses only
arithmetic. Write it in Python as soon as it needs arrays, a transform, a
file, or a plot.
\end{ruleofthumb}

There is a second, quieter reason to prefer Python once the algorithm is more
than trivial: the model becomes the specification. \file{cordic\_model.py}
is thirty lines that anybody in the group can read, run and modify, including
the person who has not yet learned \SV{}. When the architect changes the
number of guard bits, the change happens in one file, the model is rerun in a
second, the new error figure is printed, and only then does anybody touch
\RTL{}. A model buried inside a testbench package cannot play that role,
because running it requires a simulator.

\begin{tipbox}{Make the model runnable on its own}
Give every Python model a \code{\_\_main\_\_} block that prints its
parameters and its worst-case error. \file{cordic\_model.py} does exactly
that, and \code{python3 cordic\_model.py} is the fastest sanity check in the
whole flow -- no simulator, no build, one second. When somebody reports that
the hardware is wrong, run the model first.
\end{tipbox}

\begin{notebox}{Tool support for real-valued constant functions}
Constant functions that call real-valued system functions are legal in IEEE
1800-2017 and IEEE 1076-2008, but support is uneven. \tool{Verilator} 5.020
elaborates the \code{\$atan} table above correctly. \tool{Icarus Verilog}
12.0 does not: it cannot evaluate a constant function that returns an
unpacked array, and it does not accept an unpacked array \kw{localparam} at
all. Synthesis tools vary more widely still; several accept
\code{ieee.math\_real} in a constant expression but reject the same
expression inside a function used to initialise a ROM. This is the single
strongest practical argument for generating such tables from Python:
generated integer literals are accepted by everything.
\end{notebox}

\section{Four ways to couple Python to a testbench}
\label{sec:python:taxonomy}

There are exactly four places where the two worlds can meet: in a file, at
compile time, inside the simulator process, or across a process boundary.
\Cref{fig:python:coupling} shows all four.

\begin{figure}[htbp]
  \centering
  \begin{tikzpicture}[
      font=\small\sffamily,
      py/.style={draw=codekey,fill=notebg,rounded corners=2pt,
                 minimum width=21mm,minimum height=8mm,align=center},
      hdl/.style={draw=accentalt,fill=svbg,rounded corners=2pt,
                 minimum width=21mm,minimum height=8mm,align=center},
      art/.style={draw=codeframe,fill=codebg,minimum width=17mm,
                 minimum height=7mm,align=center,font=\scriptsize\sffamily},
      ar/.style={-{Latex[length=2mm]},thick,draw=codenumber},
      lbl/.style={font=\scriptsize\bfseries\sffamily,color=accent}]

    % (a) offline
    \node[lbl] at (0,1.5) {(a) offline vectors};
    \node[py]  (pa) at (0,0.7) {Python};
    \node[art] (fa) at (0,-0.2) {\file{.txt} file};
    \node[hdl] (ha) at (0,-1.1) {simulator};
    \draw[ar] (pa) -- (fa);
    \draw[ar] (fa) -- (ha);

    % (b) generated source
    \node[lbl] at (4.4,1.5) {(b) generated source};
    \node[py]  (pb) at (4.4,0.7) {Python};
    \node[art] (fb) at (4.4,-0.2) {\file{.sv} / \file{.h}};
    \node[hdl] (hb) at (4.4,-1.1) {compiler};
    \draw[ar] (pb) -- (fb);
    \draw[ar] (fb) -- (hb);

    % (c) in-process
    \node[lbl] at (8.8,1.5) {(c) in-process};
    \node[py]  (pc) at (8.8,0.7) {Python};
    \node[hdl] (hc) at (8.8,-1.1) {simulator};
    \draw[ar] (pc.south west) ++(2mm,0) -- ++(0,-1.0)
      node[midway,left,font=\scriptsize] {call};
    \draw[ar] (hc.north east) ++(-2mm,0) -- ++(0,1.0)
      node[midway,right,font=\scriptsize] {call};
    \node[art,draw=none,fill=none] at (8.8,-0.2) {one process};

    % (d) out-of-process
    \node[lbl] at (13.2,1.5) {(d) co-simulation};
    \node[py]  (pd) at (13.2,0.7) {Python server};
    \node[hdl] (hd) at (13.2,-1.1) {simulator};
    \draw[ar] (pd) -- (hd);
    \draw[ar] (hd.north west) ++(3mm,0) -- ++(0,1.0);
    \node[art,draw=none,fill=none] at (13.2,-0.2)
      {\scriptsize socket};

    \begin{scope}[on background layer]
      \node[draw=codeframe,dashed,rounded corners=3pt,
            fit=(pc)(hc),inner sep=3.5mm] {};
    \end{scope}
  \end{tikzpicture}
  \caption{The four ways to couple a Python model to a simulation. They
  differ in \emph{when} the coupling happens: (a) before the run, (b) before
  the build, (c) during the run in one process, (d) during the run across a
  process boundary.}
  \label{fig:python:coupling}
\end{figure}

\index{coupling!offline vectors}\index{coupling!generated source}
\index{coupling!in-process}\index{coupling!co-simulation}
\Cref{tab:python:tradeoffs} is the summary you should keep. The column that
usually decides is the last one: can the Python side react to what the \DUT
does, cycle by cycle? Styles (a) and (b) cannot. That is not a defect; it is
the reason they are cheap.

\begin{table}[htbp]
\centering
\caption{Trade-offs of the four coupling styles.}
\label{tab:python:tradeoffs}
\small
\begin{tabularx}{\linewidth}{@{}lXXXX@{}}
\toprule
 & \textbf{(a) offline} & \textbf{(b) generated} & \textbf{(c) in-process}
 & \textbf{(d) sockets} \\
\midrule
Run-time cost      & none & none & \(\mu\)s per call & \(\mu\)s to ms per call \\
Debuggability      & excellent: the file is the evidence & excellent: the
                     output is source & mixed: two debuggers & poor \\
Reproducibility    & excellent, if the file is archived & excellent & good,
                     if versions are pinned & fragile \\
Setup complexity   & lowest & low & high (DPI) or medium (cocotb) & high \\
Simulator portability & total & total & DPI: good; VPI: good; cocotb: good
                     & total \\
Reacts cycle by cycle & no & no & yes & yes, slowly \\
\bottomrule
\end{tabularx}
\end{table}

\section{Offline vectors}
\label{sec:python:vectors}

This is the workhorse. Master it first: it works with every simulator you
will ever meet, it needs no libraries, it survives a change of tool vendor,
and the intermediate file is a debugging artefact you can read with your
eyes.

The flow has three steps. Python generates stimulus and expected responses
into a file. The simulation reads the file, drives the \DUT and writes what
it observed into a second file. Python reads the second file and decides
whether the design passed.

\subsection{The generator}

\index{stimulus generation}
\Cref{lst:python:gen} is the generator used throughout this chapter. It
imports the golden model from \file{cordic\_model.py} -- the same module the
cocotb testbench will import later, which is the point -- and writes two
files.

\begin{pycode}[caption={The vector generator (extract from
\file{gen\_vectors.py}).},label={lst:python:gen}]
import argparse, math, random, subprocess, sys
from pathlib import Path
sys.path.insert(0, str(Path(__file__).resolve().parent))
import cordic_model as m

def hexw(value, width):
    """Two's complement hex, `width` bits, no sign, no 0x prefix."""
    return f"{value & ((1 << width) - 1):0{(width + 3) // 4}x}"

def corner_cases():
    """Directed vectors: the angles a random generator will not hit."""
    angs = [0.0, math.pi/2, -math.pi/2, math.pi, -math.pi,
            math.pi/4, -math.pi/4, 1e-4, -1e-4]
    out = [(m.X0_UNIT, 0, m.float_to_q(a, m.QA, m.AW)) for a in angs]
    return out + [(32767, 0, 0), (-32768, 0, 0), (0, 0, 0)]

def random_cases(n, rng):
    lim = m.float_to_q(math.pi, m.QA, m.AW)
    return [(rng.randint(-6000, 6000), rng.randint(-6000, 6000),
             rng.randint(-lim, lim)) for _ in range(n)]

def main():
    args = parse_args()                 # --seed, --count, --outdir
    rng = random.Random(args.seed)
    cases = corner_cases() + random_cases(args.count, rng)

    hdr = [f"seed={args.seed} count={len(cases)}",
           f"DW={m.DW} QF={m.QF} AW={m.AW} QA={m.QA} N={m.N} "
           f"GL={m.GL} GM={m.GM}",
           f"model=cordic_model.py rev={model_version()} "
           f"python={sys.version.split()[0]}",
           "columns: x0 y0 theta xexp yexp  (hex, 2's complement, 16 bit)"]

    out = Path(args.outdir); out.mkdir(parents=True, exist_ok=True)
    ftxt = open(out / "cordic_vec.txt", "w")
    fhex = open(out / "cordic_vec.hex", "w")
    for line in hdr:
        ftxt.write("// " + line + "\n")
        fhex.write("// " + line + "\n")
    for (x0, y0, th) in cases:
        xe, ye = m.cordic_rot(x0, y0, th)
        cols = [hexw(v, 16) for v in (x0, y0, th, xe, ye)]
        ftxt.write(" ".join(cols) + "\n")
        fhex.write("".join(cols) + "\n")   # one 80-bit word per line
    ftxt.close(); fhex.close()
    # $readmemh cannot report how many entries it filled: say so explicitly.
    (out / "cordic_vec.n").write_text(f"{len(cases)}\n")
\end{pycode}

Running it produces this:

\begin{shcode}
$ python3 py/gen_vectors.py --seed 20250905 --count 200 --outdir vec
\end{shcode}

\begin{txtcode}
wrote 212 vectors to vec/cordic_vec.{txt,hex}

$ head -6 vec/cordic_vec.txt
// seed=20250905 count=212
// DW=16 QF=14 AW=16 QA=13 N=14 GL=2 GM=2
// model=cordic_model.py rev=no-git python=3.11.15
// columns: x0 y0 theta xexp yexp   (hex, two's complement, 16 bits)
26dd 0000 0000 3fff 0001
26dd 0000 3244 0001 3fff
\end{txtcode}

\subsection{Deciding the file format}

Four decisions, and each of them will bite somebody if you get it wrong.

\index{vector file!format}
\textbf{Hex or decimal?} Hex. A hex field has a fixed width, which makes the
file column-aligned and diffable, and it maps one to one onto the bit pattern
on the wire. Decimal forces you to decide, in the file, whether the value is
signed -- and then to make the reader agree. Write hex, and treat the file as
a picture of the pins.

\textbf{One transaction per line.} Never split a transaction across lines and
never put two on one line. Line number equals transaction number is the most
useful invariant a vector file can have: when the simulation reports a
mismatch on transaction 137 you open the file at line 141 (137 plus the
four-line header) and look at it.

\textbf{A header.} The first lines must record the seed, the parameter set
and the model revision that produced the file. A vector file without a header
is a liability: six months later nobody can tell whether it was generated for
\code{N = 14} or \code{N = 16}, and a testbench that reads it will happily
report thousands of mismatches that mean nothing.

\textbf{A comment convention.} Use \code{//}. It is what \code{\$readmemh}
accepts, it is what a C++ \code{ifstream} loop can skip with two characters
of test, and \code{std.textio} can skip it with one \code{read} of a
character. A \code{\#} convention would be more Pythonic and would break
\code{\$readmemh}.

\subsection{Stimulus only, or stimulus and expected response?}

The generator above writes both: the three inputs and the two expected
outputs, on the same line. That is the right default, and it is worth being
clear about why, because the alternative is common and defensible.

If the file contains only stimulus, the checking has to happen somewhere
else -- either in the HDL testbench, which then needs a model, or in Python
afterwards, which then needs the results file and the stimulus file to be
correlated by line number. If the file contains both, the HDL testbench is a
pure comparator: it drives, it samples, it compares two integers. It needs no
model, no real arithmetic and no knowledge of what a CORDIC is. That is a
testbench a second-year student can maintain, and it is the same testbench
whatever the design computes.

The cost is that the file is bigger and that it is now tied to one version of
the model. Regenerate it whenever the model changes -- which is what the Make
rule in \cref{sec:python:vectors} is for.

Two refinements are worth knowing. First, when the response is large
(an image, a block of samples) put it in a separate file and keep the index
in the stimulus file; a vector file you cannot open in an editor has lost
most of its value. Second, when the transaction has optional fields, put a
one-character command at the start of the line -- \code{W} for write,
\code{R} for read, \code{\#} would clash with the comment convention -- and
dispatch on it in the reader. That turns the vector file into a small
instruction stream, and it is how most directed-test frameworks start.

\begin{tipbox}{The vector file header}
Put the seed, every parameter the model depends on, the model file's revision
and the Python version in the header, each on its own \code{//} line. Then
make the testbench print the header it read at time zero. A regression log
that contains the header of the vector file it consumed is self-documenting.
\end{tipbox}

\subsection{Reading it in \SV{}: \code{\$readmemh}}

\index{\$readmemh}
\code{\$readmemh} loads whitespace-separated hex tokens into successive
entries of an array. It is the fastest thing to write and the easiest thing
to get subtly wrong. The version below packs the five 16-bit fields of one
transaction into a single 80-bit word, so that one array entry is one
transaction.

\begin{svcode}[caption={Loading vectors with \code{\$readmemh}
(\file{tb\_vec\_readmem.sv}).},label={lst:python:readmem}]
module tb_vec_readmem;
  import cordic_pkg::*;

  localparam int MAXV = 4096;
  localparam int VW   = 5*DW;          // 80 bits: x0 y0 theta xexp yexp

  logic [VW-1:0] vec [0:MAXV-1];
  int            nvec;

  // How many entries did $readmemh actually fill?  It will not tell us, so
  // the generator writes the count into a companion file.
  function automatic int read_count(input string name);
    int fd, n;
    fd = $fopen(name, "r");
    if (fd == 0) $fatal(1, "cannot open %s", name);
    if ($fscanf(fd, "%d", n) != 1) $fatal(1, "bad count file %s", name);
    $fclose(fd);
    return n;
  endfunction

  initial begin
    if (!$value$plusargs("VEC=%s", vecfile)) vecfile = "vec/cordic_vec.hex";
    nvec = read_count({vecfile.substr(0, vecfile.len()-5), ".n"});
    $readmemh(vecfile, vec, 0, nvec-1);
    $display("[tb] loaded %0d vectors from %s", nvec, vecfile);
    ...
    for (int i = 0; i < nvec; i++) begin
      data_t  x0, y0, xe, ye;
      angle_t th;
      {x0, y0, th, xe, ye} = vec[i];      // unpack one 80-bit word

      @(negedge clk);
      req_x = x0; req_y = y0; req_theta = th; req_valid = 1'b1;
      do @(posedge clk); while (!req_ready);
      @(negedge clk); req_valid = 1'b0;
      do @(posedge clk); while (!rsp_valid);

      if (rsp_x !== xe || rsp_y !== ye) errors++;
    end
    $display("[tb] %0d vectors, %0d mismatches", nvec, errors);
    $display("[tb] %s", errors == 0 ? "PASS" : "FAIL");
    $finish;
  end
endmodule
\end{svcode}

Built and run with \tool{Verilator}:

\begin{shcode}
$ verilator --binary --timing --top-module tb_vec_readmem \
      sv/cordic_pkg.sv sv/cordic_rot.sv sv/tb_vec_readmem.sv \
      -o Vtb_readmem --Mdir obj_readmem
$ ./obj_readmem/Vtb_readmem
\end{shcode}

\begin{txtcode}
[tb] loaded 212 vectors from vec/cordic_vec.hex
[tb] 212 vectors, 0 mismatches
[tb] PASS
- sv/tb_vec_readmem.sv:73: Verilog $finish
\end{txtcode}

Note the unpacking line, \code{\{x0, y0, th, xe, ye\} = vec[i];}. A
concatenation on the left-hand side is a \SV{} idiom with no \VHDL{}
equivalent: it splits a vector across several targets by width, left to
right. It is the natural partner of a packed vector file.

\begin{gotcha}{\code{\$readmemh} reads a short file in silence}
The first version of this testbench filled the array with \code{'x} and then
counted how many entries \code{\$readmemh} had overwritten. Under
\tool{Verilator} that count came back as 4096 -- the size of the array, not
the size of the file -- because \tool{Verilator} is a two-state simulator and
the \code{'x} initialisation had become zero. The run then drove 3884
all-zero transactions, which the \DUT dutifully turned into all-zero
responses, and reported:

\begin{txtcode}
[tb] loaded 4096 vectors from vec/cordic_vec.hex
[tb] 4096 vectors, 0 mismatches
[tb] PASS
\end{txtcode}

A green result from a testbench that ran almost nothing. \code{\$readmemh}
does not report how many entries it filled, and if the file is shorter than
the array the remaining entries simply keep their old value. In a four-state
simulator they stay \code{'x} and you will usually notice; under
\tool{Verilator} they are zero and you will not. The fix is the one in
\cref{lst:python:readmem}: have the generator write the count into a
companion file and pass an explicit address range,
\code{\$readmemh(file, vec, 0, nvec-1)}.
\end{gotcha}

Two more traps in the same system task. First, the address arguments are
\emph{addresses}, not counts, and they are inclusive: \code{\$readmemh(f, m,
0, 9)} loads ten entries. Second, the file may contain \code{@} directives --
a line \code{@1F} sets the load address to \code{'h1F} -- and a stray
\code{@} in a comment your generator wrote will silently relocate everything
that follows. Generate the file, never hand-edit it, and never let a
free-text field into it.

\subsection{Reading it line by line}

\index{\$fgets}\index{\$sscanf}\index{\$feof}
When the file has structure -- several columns, optional fields, a per-line
command -- read it a line at a time. This is also the version to reach for
when you want the file to stream: \code{\$readmemh} loads everything into
memory before the simulation starts, which is fine for a thousand vectors and
not fine for ten million.

\begin{svcode}[caption={Line-by-line reading with
\code{\$fgets}/\code{\$sscanf} (\file{tb\_vec\_file.sv}).},%
label={lst:python:fgets}]
  int    fd_in, fd_out, code, nline = 0, nvec = 0, errors = 0;
  string line, dir;

  initial begin
    if (!$value$plusargs("DIR=%s", dir)) dir = "vec";
    fd_in  = $fopen({dir, "/cordic_vec.txt"}, "r");
    if (fd_in == 0) $fatal(1, "cannot open the vector file");
    fd_out = $fopen({dir, "/cordic_res.txt"}, "w");
    ...
    while (!$feof(fd_in)) begin
      logic [15:0] x0, y0, th, xe, ye;
      code = $fgets(line, fd_in);
      nline++;
      if (code <= 0) break;                       // short read or EOF
      if (line.len() < 2 || line.substr(0,1) == "//") continue;
      code = $sscanf(line, "%h %h %h %h %h", x0, y0, th, xe, ye);
      if (code != 5) begin
        $display("[tb] line %0d: expected 5 fields, got %0d -- skipped",
                 nline, code);
        continue;
      end
      nvec++;
      ... drive the DUT, wait for the response ...
      $fdisplay(fd_out, "%04h %04h %04h %04h %04h",
                x0, y0, th, rsp_x, rsp_y);
      if (rsp_x !== xe || rsp_y !== ye) errors++;
    end
    $fclose(fd_in); $fclose(fd_out);
    $display("[tb] %0d lines, %0d vectors, %0d mismatches",
             nline, nvec, errors);
  end
\end{svcode}

\begin{txtcode}
$ ./obj_file/Vtb_file +DIR=vec
[tb] 217 lines, 212 vectors, 0 mismatches
[tb] PASS
\end{txtcode}

\code{\$sscanf} returns the number of fields it successfully converted, and
that return value is the only line-format check you get. Test it. A file
whose last line is truncated because the generator was interrupted will
otherwise be read as a valid transaction with three of five fields left at
their previous values.

\begin{gotcha}{\code{\$feof} is not a look-ahead}
\code{\$feof} returns non-zero only \emph{after} a read has hit the end of
the file, exactly as in C. A \code{while (!\$feof(fd))} loop therefore
executes one extra iteration on an empty last read. That is why
\cref{lst:python:fgets} also checks the return value of \code{\$fgets} and
breaks. \VHDL{}'s \code{endfile} is a true look-ahead and does not have this
problem, which is precisely why a \VHDL{} designer forgets to add the check.
\end{gotcha}

\subsection{The same file in \VHDL{}}

\index{textio}\index{hread}
For comparison, and because you will still be writing \VHDL{} testbenches for
some time, here is the same loop with \code{std.textio}. \VHDL{}-2008's
\code{hread} reads a hex value into a \code{std\_logic\_vector} of matching
width; before 2008 you needed \code{ieee.std\_logic\_textio} from a vendor
library.

\begin{vhdlcode}[caption={The same vector file read with \code{std.textio}
(\file{tb\_vec\_vhdl.vhd}).},label={lst:python:textio}]
  stim : process
    file     f      : text;
    variable l      : line;
    variable status : file_open_status;
    variable c      : character;
    variable ok     : boolean;
    variable x0, y0, xe, ye : std_logic_vector(DW-1 downto 0);
    variable th             : std_logic_vector(AW-1 downto 0);
    variable nvec, errors   : natural := 0;
  begin
    file_open(status, f, VEC, read_mode);
    assert status = open_ok report "cannot open " & VEC severity failure;

    wait for 33 ns;
    rst_n <= '1';

    while not endfile(f) loop
      readline(f, l);
      next when l'length = 0;
      read(l, c, ok);                       -- peek at the first character
      next when not ok or c = '/';          -- skip '//' comments
      l := new string'(c & l.all);          -- put the character back

      hread(l, x0, ok); next when not ok;
      hread(l, y0);
      hread(l, th);
      hread(l, xe);
      hread(l, ye);
      nvec := nvec + 1;

      drive_and_wait(x0, y0, th);          -- the usual handshake

      if std_logic_vector(rsp_x) /= xe or std_logic_vector(rsp_y) /= ye then
        errors := errors + 1;
      end if;
    end loop;
    file_close(f);
    report "vectors = " & integer'image(nvec) &
           "  mismatches = " & integer'image(errors);
    assert errors = 0 report "FAIL" severity failure;
    report "PASS";
    wait;
  end process;
\end{vhdlcode}

\begin{shcode}
$ vlib work
$ vcom -2008 vhdl/cordic_pkg.vhd vhdl/cordic_rot.vhd vhdl/tb_vec_vhdl.vhd
$ vsim -c -do "run -all; quit" tb_vec_vhdl
\end{shcode}

Everything the testbench prints comes from \code{report}, so the two lines
that matter read the same whatever tool ran it, once the per-simulator prefix
is stripped off:

\begin{txtcode}
vectors = 212  mismatches = 0
PASS
\end{txtcode}

The interesting difference is not the syntax but the type discipline.
\code{hread} knows how many hex digits to consume because the target has a
width; \code{\$sscanf} does not, and will happily read six digits into a
16-bit variable and throw the top away. \VHDL{} makes the format explicit and
\SV{} makes it your problem -- the same trade you meet everywhere else
between the two languages.

\subsection{The same file in a \tool{Verilator} C++ testbench}

\index{Verilator!C++ testbench}
If your testbench is C++ (\cref{ch:verilator}), the reader is an
\code{ifstream} loop and the file format does not change at all. That is the
whole point of the style: one file, four readers.

\begin{cppcode}[caption={The vector file in a \tool{Verilator} C++
testbench.},label={lst:python:cppread}]
struct Vec { int16_t x0, y0, th, xe, ye; };

static int16_t hex16(const std::string& s) {
  return static_cast<int16_t>(std::stoul(s, nullptr, 16));
}

static std::vector<Vec> load(const char* path) {
  std::vector<Vec> v;
  std::ifstream f(path);
  if (!f) { std::fprintf(stderr, "cannot open %s\n", path); std::exit(2); }
  std::string line;
  while (std::getline(f, line)) {
    if (line.size() >= 2 && line[0] == '/' && line[1] == '/') continue;
    std::istringstream is(line);
    std::string a, b, c, d, e;
    if (!(is >> a >> b >> c >> d >> e)) continue;
    v.push_back({hex16(a), hex16(b), hex16(c), hex16(d), hex16(e)});
  }
  return v;
}

  // The driving loop is the same handshake as before: assert req_valid_i,
  // tick until req_ready_o, tick once more, then tick until rsp_valid_o.
  for (const Vec& v : vecs) {
    dut->req_x_i = v.x0; dut->req_y_i = v.y0; dut->req_theta_i = v.th;
    dut->req_valid_i = 1;
    dut->eval();
    while (!dut->req_ready_o) tick();   // stall until the DUT accepts
    tick();                            // the accepting clock edge
    dut->req_valid_i = 0;
    dut->eval();
    while (!dut->rsp_valid_o) tick();
    if (static_cast<int16_t>(dut->rsp_x_o) != v.xe ||
        static_cast<int16_t>(dut->rsp_y_o) != v.ye) ++errors;
  }
\end{cppcode}

\begin{txtcode}
$ ./obj_cpp/Vtb_cpp vec/cordic_vec.txt
[cpp] 212 vectors, 0 mismatches, 6788 half-cycles
[cpp] PASS
\end{txtcode}

\begin{gotcha}{Signed values in a hex vector file}
\code{26dd 0000 cdbc 0001 c000} contains three negative numbers, and nothing
in the file says so. Every reader must be told the width and must sign-extend
itself. In C++ the cast is \code{static\_cast<int16\_t>(...)}; in Python it is
an explicit \code{v - 65536 if v \& 0x8000 else v}; in \SV{} it is
\code{\$signed(...)} or an assignment to a \kw{logic} \kw{signed} variable;
in \VHDL{} it is the \code{signed(...)} type conversion. Get it wrong on
\emph{one} side and the mismatch report will show you \code{52668} where you
expected \code{-12868}, which at least is obvious. Get it wrong on
\emph{both} sides in the same way and the comparison still passes -- and your
tolerance check on the floating point reference will be the only thing that
catches it. Always print at least one negative expected value at time zero
and check it by eye.
\end{gotcha}

\subsection{Checking the other direction}

\index{post-processing}
The simulation wrote \file{cordic\_res.txt}. Python now reads it, compares
against the \emph{floating point} reference (not the bit-exact one -- the
bit-exact comparison already happened inside the simulation) and reports the
error in LSBs.

\begin{pycode}[caption={Post-processing and plotting
(\file{check\_results.py}).},label={lst:python:check}]
import numpy as np
import matplotlib
matplotlib.use("Agg")            # no display in CI
import matplotlib.pyplot as plt
import cordic_model as m

def s16(v):
    return v - 65536 if v & 0x8000 else v

def load(path):
    rows = []
    for line in Path(path).read_text().splitlines():
        line = line.split("//")[0].strip()
        if not line:
            continue
        rows.append([s16(int(t, 16)) for t in line.split()])
    return np.array(rows, dtype=np.int64)

def main():
    r = load(args.results)
    x0, y0, th, xd, yd = (r[:, i] for i in range(5))

    theta = th / float(1 << m.QA)
    xf = m.K * (x0 * np.cos(theta) - y0 * np.sin(theta))
    yf = m.K * (x0 * np.sin(theta) + y0 * np.cos(theta))

    ex, ey = xd - xf, yd - yf            # error in LSB of the output format
    err = np.maximum(np.abs(ex), np.abs(ey))

    # The float reference does not saturate; the DUT does.  Exclude the
    # vectors that hit a rail, and say how many they were.
    hi, lo = (1 << (m.DW - 1)) - 1, -(1 << (m.DW - 1))
    sat = (xd >= hi) | (xd <= lo) | (yd >= hi) | (yd <= lo)
    err = np.where(sat, 0.0, err)

    print(f"vectors        : {len(r)}")
    print(f"saturated      : {int(sat.sum())} (excluded)")
    print(f"worst |error|  : {err.max():.3f} LSB")
    print(f"tolerance      : {args.tol_lsb} LSB")
    bad = np.flatnonzero(err > args.tol_lsb)
    print(f"outside tol    : {len(bad)}")

    if args.plot:
        fig, ax = plt.subplots(figsize=(7, 3.2))
        ax.plot(theta[~sat], ex[~sat], lw=0.8, label="x error")
        ax.plot(theta[~sat], ey[~sat], lw=0.8, label="y error")
        for s in (args.tol_lsb, -args.tol_lsb):
            ax.axhline(s, ls="--", lw=0.7, color="0.4")
        ax.set_xlabel("theta [rad]")
        ax.set_ylabel("error [LSB of Q1.14]")
        ax.legend(loc="upper right", fontsize=8); ax.grid(alpha=0.3)
        fig.tight_layout(); fig.savefig(args.plot, dpi=150)

    return 1 if len(bad) else 0
\end{pycode}

Run over a 721-point sweep of the unit circle:

\begin{txtcode}
$ python3 py/check_results.py sweep/cordic_res.txt --plot sweep/err.png
vectors        : 721
saturated      : 0 (excluded)
worst |error|  : 3.597 LSB
rms   error    : 1.115 LSB
tolerance      : 12.0 LSB
outside tol    : 0
plot           : sweep/err.png
\end{txtcode}

\placeholder{6cm}{The error-versus-angle plot produced by
\file{check\_results.py}: two traces, x error and y error in LSBs of Q1.14,
against theta from -pi to +pi. Both stay inside a band of about
\(\pm 3.6\) LSB, with the largest excursions near \(\pm \pi/2\) and
\(\pm \pi\) where the coarse rotation stage runs and the residual angle is
largest. The dashed lines at \(\pm 12\) LSB are the acceptance tolerance;
the traces never approach them.}

\index{tolerance}
That number, 3.597 LSB, is the whole reason to keep the floating point
reference around. The bit-exact model told you the \RTL{} matches the model.
The floating point reference tells you the \emph{model} is a good enough
approximation of a rotation -- which is a different question, and the one
your supervisor will actually ask.

\begin{gotcha}{A floating point golden model cannot be compared exactly}
A bit-exact Python model and the \RTL{} must agree exactly: any difference is
a bug in one of them, and the check is \code{==}. A floating point reference
and a fixed point \DUT{} will \emph{never} agree exactly, and comparing them
needs a tolerance. Choosing that tolerance is a design decision, not a
testbench detail:
\begin{itemize}
\item Derive it. For an \(N\)-step CORDIC the angle quantisation contributes
      about \(2^{-N}\) radian and the datapath truncation about one LSB per
      stage in the worst case; for \(N = 14\) with two guard bits the bound is
      a little under 5 LSB. The testbench in \refverificationbook uses 12 LSB, which
      leaves margin without hiding a real bug.
\item Never tune it upward until the test passes. If you must, write down in
      the commit message why the bound moved.
\item Exclude the cases the reference cannot model. The float reference does
      not saturate; the \DUT{} does. Before that exclusion was added, two
      directed vectors with \code{x0 = 32767} produced a reported worst error
      of \(21193\) LSB and a red build. The vectors were fine; the comparison
      was wrong.
\end{itemize}
\end{gotcha}

\subsection{Reproducibility}

\index{reproducibility}\index{seed}
A vector file is data, and data in a verification flow has to be
reproducible. Four rules.

\textbf{The seed is an input, not a constant.} Pass it on the command line,
default it to something fixed, print it, and put it in the file header. A
regression that finds a bug with \code{--seed 4711} must be repeatable by
typing \code{--seed 4711}.

\textbf{Pin the versions.} \code{random.Random(seed)} in CPython is stable
across versions in practice, but \code{numpy}'s legacy \code{np.random.seed}
is explicitly not guaranteed to be, and \code{set} iteration order is not
either. If a vector file must be byte-identical next year, use
\code{random.Random}, sort anything that came out of a set or a dict, and
record the Python version in the header.

\textbf{Store the file, or store the generator -- decide.} Committing the
vector file makes the repository self-contained and the history readable, at
the cost of a binary-ish blob in every diff. Regenerating it keeps the
repository small but means a broken generator breaks every historical commit.
The usual compromise: commit a small directed set (a few hundred vectors,
reviewable) and regenerate the large random set in CI.

\textbf{Make the build depend on the generator.} This is the one people
forget.

\begin{makecode}
# The vector file depends on the generator AND on the model: change either
# one and the vectors are rebuilt.
vec/cordic_vec.txt vec/cordic_vec.hex vec/cordic_vec.n: \
        py/gen_vectors.py py/cordic_model.py
	$(PYTHON) py/gen_vectors.py --seed $(SEED) --count $(COUNT) \
	    --outdir vec
\end{makecode}

\begin{ruleofthumb}
Every generated artefact must have a Make rule whose prerequisites include
both the generator and everything the generator imports. Otherwise the day
you fix the model is the day your regression starts lying to you.
\end{ruleofthumb}

\section{Generated source}
\label{sec:python:generated}

\index{code generation}
The second style moves the coupling one step earlier: Python writes source
code, and the compiler consumes it. The advantage over offline vectors is
that the result becomes a constant in the compiled design, so there is no
run-time file access and no timing at all. The cost is a build step and a
generated file that can go stale.

\subsection{One generator, two outputs}

\index{jinja2}
\Cref{lst:python:gentbl} emits the CORDIC constants both as a C++ header for
the \tool{Verilator} testbench and as a \SV{} package for the \RTL{}. It uses
\code{jinja2}, which is worth the dependency as soon as the template has more
than one loop in it: the alternative is a function full of string
concatenation that nobody can read.

\begin{pycode}[caption={Emitting a C++ header and a \SV{} package from one
model (\file{gen\_tables.py}).},label={lst:python:gentbl}]
from datetime import date
from jinja2 import Template
import cordic_model as m

BANNER = ("// GENERATED by gen_tables.py on {date} -- DO NOT EDIT.\n"
          "// Source of truth: cordic_model.py (N={N}, IQA={IQA}).\n")

CPP = Template("""{{banner}}
#pragma once
#include <array>
#include <cstdint>

namespace cordic_ref {

constexpr int    kN = {{N}}, kDW = {{DW}}, kQF = {{QF}};
constexpr int    kAW = {{AW}}, kQA = {{QA}}, kX0Unit = {{X0}};
constexpr double kGain = {{K}};

constexpr std::array<int32_t, kN> kAtan = {
{%- for v in atan %}
    {{ v }}{{ "," if not loop.last }}
{%- endfor %}
};

}  // namespace cordic_ref
""")

SV = Template("""{{banner}}
package cordic_tbl_pkg;
  localparam int unsigned N   = {{N}};
  localparam int unsigned IQA = {{IQA}};
  localparam int signed   X0_UNIT = {{X0}};

  localparam int signed ATAN_ROM [{{N}}] = '{
{%- for v in atan %}
    {{ v }}{{ "," if not loop.last }}
{%- endfor %}
  };
endpackage : cordic_tbl_pkg
""")

def main():
    out = Path(sys.argv[1] if len(sys.argv) > 1 else "gen")
    out.mkdir(parents=True, exist_ok=True)
    ctx = dict(banner=BANNER.format(date=date.today(), N=m.N, IQA=m.IQA),
               N=m.N, DW=m.DW, QF=m.QF, AW=m.AW, QA=m.QA, IQA=m.IQA,
               K=repr(m.K), X0=m.X0_UNIT, atan=m.ATAN)
    (out / "cordic_tbl.h").write_text(CPP.render(**ctx))
    (out / "cordic_tbl_pkg.sv").write_text(SV.render(**ctx))
\end{pycode}

The output is boring, which is the point:

\begin{txtcode}
$ python3 py/gen_tables.py gen && head -20 gen/cordic_tbl.h
// GENERATED by gen_tables.py on 2026-09-05 -- DO NOT EDIT.
// Source of truth: cordic_model.py (N=14, IQA=15).

#pragma once
#include <array>
#include <cstdint>

namespace cordic_ref {

constexpr int    kN      = 14;
constexpr int    kDW     = 16, kQF = 14;
constexpr int    kAW     = 16, kQA = 13;
constexpr double kGain   = 1.6467602540312922;
constexpr int    kX0Unit = 9949;

constexpr std::array<int32_t, kN> kAtan = {
    25736,
    15193,
    8027,
...
\end{txtcode}

\subsection{A package, not an include}

\index{package!generated}\index{\`{}include}
You can emit either a \SV{} package or a \file{.svh} file that gets
\code{`include}d into a module. Emit the package.

An \code{`include} is textual substitution done by the preprocessor. Include
it in two modules and you get two copies of the constants, with no
guarantee -- and no check -- that both compiled from the same generation.
Include it twice into the same scope by accident and you get a duplicate
declaration error whose message points at a line number in a file you did not
write. A package is a named scope with a single elaboration: import it in ten
places and there is still exactly one \code{ATAN\_ROM}, and the tool will
tell you if two versions of the package are in the compilation library. This
is the same argument you already accept in \VHDL{}, where nobody would dream
of using textual inclusion instead of a package.

The one case for an \code{`include} is a fragment that is not a legal
compilation unit on its own -- a list of port connections, a block of
declarations that must land inside a module. If you find yourself generating
those, consider whether the thing you actually want is a parameterised module
or an interface.

\subsection{The arctangent ROM: generate or elaborate?}

\index{ROM!arctangent}
\file{cordic\_pkg.sv} computes the arctangent table at elaboration with
\code{\$atan}, as shown at the start of this chapter. \file{gen\_tables.py}
computes the same table in Python and writes integer literals. Both are
correct. Which one you should use depends entirely on your tools.

\tool{Verilator} 5.020 handles the elaborated version:

\begin{txtcode}
$ ./obj_atan/Vtb_atan
ATAN_ROM = 25736 15193 8027 4075 2045 1024 512 256 128 64 32 16 8 4
X0_UNIT = 9949, HALF_PI = 51472
\end{txtcode}

which is exactly the Python table:

\begin{txtcode}
$ python3 py/cordic_model.py
N = 14, K = 1.646760, 1/K = 0.607253, X0_UNIT = 9949
HALF_PI = 51472, PI = 102944
atan table: [25736, 15193, 8027, 4075, 2045, 1024, 512, 256, 128, 64, 32,
             16, 8, 4]
worst |error| over [-pi, pi] = 0.000233  (3.82 LSB)
\end{txtcode}

\tool{Icarus Verilog} 12.0 does not:

\begin{txtcode}
$ iverilog -g2012 -DSYNTHESIS -o /dev/null sv/cordic_pkg.sv sv/cordic_rot.sv
sv/cordic_pkg.sv:106: error: Array rom needs an array index here.
sv/cordic_pkg.sv:109: error: Unable to evaluate parameter ATAN_ROM value:
                             cordic_pkg.build_atan_rom()
2 error(s) during elaboration.
\end{txtcode}

That is the deciding factor, and it is a mundane one. Keep the constant
function if every tool in your flow -- simulator, linter, synthesiser,
formal tool -- accepts it, because it removes a build step and keeps the
design self-contained. Generate the table the moment one of them does not.

\begin{gotcha}{The generated file that nobody regenerated}
The failure mode of style (b) is always the same. Somebody changes \code{N}
in the model, runs the simulation, sees it pass, and ships. The generated
package still holds the old table, because nothing forced it to be rebuilt.
Two defences, and you want both:
\begin{itemize}
\item a Make rule with the generator and the model as prerequisites, as in
      \cref{sec:python:vectors};
\item a CI job that regenerates into a scratch directory and diffs.
\end{itemize}
\begin{txtcode}
$ make gen-check
python3 py/gen_tables.py .gencheck
diff -r gen .gencheck || (echo "generated files are stale: run 'make gen'"; \
                          exit 1)
Files gen/cordic_tbl_pkg.sv and .gencheck/cordic_tbl_pkg.sv differ
generated files are stale: run 'make gen'
make: *** [Makefile:30: gen-check] Error 1
\end{txtcode}
Also put \code{GENERATED -- DO NOT EDIT} and the name of the generator in the
first line of every emitted file. Somebody will edit it anyway, but at least
the diff will show that they did.
\end{gotcha}

\subsection{In version control, or not?}

Both answers are defensible and the argument is old. Keeping generated files
out of the repository is cleaner in principle: one source of truth, no
possibility of a stale artefact, no noisy diffs. Keeping them in makes the
repository buildable by somebody who does not have your Python environment,
makes code review of the generated \RTL{} possible, and makes \code{git
blame} on a constant meaningful.

For a thesis or a course project, commit them, and add the \code{gen-check}
job above. Your reviewer will not have \code{jinja2}. For a large industrial
flow with a curated build environment, generate them.

\begin{notebox}{Register-map generation}
\index{IP-XACT}\index{PeakRDL}
Generating a constant table is the toy version of an idea the industry runs
at scale. A register map -- addresses, fields, access types, reset values --
is written once in a machine-readable form and everything else is generated
from it: the \RTL{} register file, the C header for firmware, the UVM
register model (\refverificationbook), the documentation, and the address decoder in
the testbench. The formats you will meet are IP-XACT (IEEE 1685, XML,
heavyweight, universally supported by commercial tools) and SystemRDL
(Accellera, far more readable). On the open-source side, \code{peakrdl} reads
SystemRDL and emits \SV{}, VHDL, C headers, HTML and a UVM model;
\code{corsair} does something similar from JSON or YAML. If your design has
more than about ten registers, learn one of these instead of maintaining four
copies of the same table by hand.
\end{notebox}

\section{In-process, part 1: DPI-C}
\label{sec:python:dpi}

\index{DPI-C}
The Direct Programming Interface is the standard way for \SV{} to call C and
for C to call \SV{}. It is defined in IEEE 1800-2017 clauses 35 and H, it is
supported by every commercial simulator and by \tool{Verilator}, and it is
the mechanism underneath most of the Python integrations you will meet.
\VHDL{} has no equivalent in the standard: VHPI (IEEE 1076-2008 clause 22) is
a procedural interface, not a foreign-function interface, and the practical
route in \VHDL{} is the vendor-specific \code{foreign} attribute, whose
spelling, calling convention and capabilities differ from one tool to the
next.

\subsection{Declaring an import}

\begin{svcode}[caption={DPI import and export declarations.},%
label={lst:python:dpidecl}]
  // Pure: same inputs give the same result, no side effects, no access to
  // the design.  The tool may cache the result or elide the call.
  import "DPI-C" pure function real c_cordic_k();

  // Context: the function needs to know which SystemVerilog scope called
  // it, because it will call back in (or use svGetScope).
  import "DPI-C" context function void c_cordic_ref(
      input  int x0, input int y0, input int theta,
      output real xr, output real yr);

  // Neither: the default.  The function may have side effects but may not
  // call back into SystemVerilog.
  import "DPI-C" function void c_show_types(
      input logic [31:0] l, input bit [31:0] b);

  // Export: make a SystemVerilog function callable from C.
  export "DPI-C" function sv_note;
  function void sv_note(input string msg);
    $display("[SV] note from C: %s", msg);
  endfunction
\end{svcode}

\index{DPI-C!pure}\index{DPI-C!context}
\kw{pure} and \kw{context} are the two qualifiers that matter. A \kw{pure}
function must depend only on its arguments, must not touch the design, must
not do I/O, and must not call back into \SV{}. In exchange, the simulator may
call it once and reuse the answer, or skip the call entirely if the result is
unused. A \kw{context} function is given the calling scope, which it needs if
it wants to call an exported task or read a signal through
\code{svGetScope}/\code{svSetScope}. Marking something \kw{pure} that is not
is a genuinely nasty bug, because the symptom is a stale value rather than a
crash.

\index{svGetScope}\index{DPI-C!open array}
Two parts of the interface deserve a sentence each because they are where the
examples in vendor documentation stop. \code{svGetScope()} returns an opaque
handle to the \SV{} scope a \kw{context} function was called from, and
\code{svSetScope(h)} makes that scope current before you call an exported
task. You need the pair when the C side must call back into a \emph{specific}
instance -- a driver in one agent rather than another -- and the usual
pattern is to have each instance register its scope in an \code{initial}
block, store the handle in a C table, and set it before each callback. Open
arrays are the other gap: an argument declared \code{input int arr []} on the
\SV{} side arrives as an \code{svOpenArrayHandle}, and you must call
\code{svSize}, \code{svLow}, \code{svHigh} and \code{svGetArrElemPtr1} to get
at the elements rather than indexing a C array. Sized arrays
(\code{input int arr [8]}) arrive as an ordinary pointer and are far simpler;
prefer them unless the size genuinely varies.

\begin{table}[htbp]
\centering
\caption{DPI type mapping (IEEE 1800-2017, Annex H). \code{svBitVecVal} and
\code{svLogicVecVal} are arrays of 32-bit chunks.}
\label{tab:python:dpitypes}
\small
\begin{tabular}{@{}lll@{}}
\toprule
\textbf{\SV{}} & \textbf{C} & \textbf{Notes} \\
\midrule
\code{byte}                & \code{char}                & 8-bit, two-state \\
\code{shortint}            & \code{short int}           & 16-bit \\
\code{int}                 & \code{int}                 & 32-bit \\
\code{longint}             & \code{long long}           & 64-bit \\
\code{real}                & \code{double}              & \\
\code{shortreal}           & \code{float}               & \\
\code{string}              & \code{const char*}         & read-only, borrowed \\
\code{chandle}             & \code{void*}               & opaque pointer \\
\code{bit}                 & \code{svBit}               & one two-state bit \\
\code{logic}, \code{reg}   & \code{svLogic}             & one four-state bit \\
\code{bit [N-1:0]}         & \code{svBitVecVal*}        & packed, 2-state \\
\code{logic [N-1:0]}       & \code{svLogicVecVal*}      & packed, 4-state:
                                                          \code{aval}+\code{bval} \\
\code{open array} \code{[]} & \code{svOpenArrayHandle}   & use the
                                                          \code{svSize}/\code{svGetArrElemPtr} API \\
\bottomrule
\end{tabular}
\end{table}

\subsection{A working example}

The C side computes the CORDIC reference in double precision and calls back
into \SV{} once, to prove that \kw{context} and \code{export} work.

\begin{cppcode}[caption={The C reference model (\file{cordic\_dpi.c}).},%
label={lst:python:dpic}]
#include <math.h>
#include <stdio.h>
#include "svdpi.h"

/* The simulator may compile this file with a C++ compiler (Verilator does),
   so every DPI symbol must have C linkage in both cases. */
#ifdef __cplusplus
extern "C" {
#endif

extern void sv_note(const char* msg);   /* exported from SystemVerilog */

#define QF 14
#define QA 13
#define NIT 14

static double cordic_gain(void) {          /* K = prod sqrt(1 + 2^-2i) */
  double k = 1.0;
  for (int i = 0; i < NIT; i++) k *= sqrt(1.0 + pow(2.0, -2.0 * i));
  return k;
}
double c_cordic_k(void) { return cordic_gain(); }

void c_cordic_ref(int x0, int y0, int theta, double* xr, double* yr) {
  static int first = 1;
  double th = (double)theta / (double)(1 << QA);
  double k  = cordic_gain();
  *xr = k * ((double)x0 * cos(th) - (double)y0 * sin(th));
  *yr = k * ((double)x0 * sin(th) + (double)y0 * cos(th));
  if (first) {
    char buf[128];
    snprintf(buf, sizeof buf, "C reference active, K = %.9f", k);
    sv_note(buf);                 /* needs `context' on the import */
    first = 0;
  }
}

/* 4-state vs 2-state: `logic' arrives as svLogicVecVal (aval + bval),
   `bit' as svBitVecVal (one word).  Getting this wrong is silent. */
void c_show_types(const svLogicVecVal* l, const svBitVecVal* b) {
  printf("[C] logic: aval=0x%08x bval=0x%08x   bit: 0x%08x\n",
         (unsigned)l->aval, (unsigned)l->bval, (unsigned)*b);
}

#ifdef __cplusplus
}
#endif
\end{cppcode}

\tool{Verilator} compiles the C file as part of the build:

\begin{shcode}
$ verilator --binary --timing --top-module tb_dpi \
      sv/cordic_pkg.sv sv/tb_dpi.sv cpp/cordic_dpi.cpp \
      -o Vtb_dpi --Mdir obj_dpi
$ ./obj_dpi/Vtb_dpi
\end{shcode}

\begin{txtcode}
[SV] K from C = 1.646760254
[C] logic: aval=0xdeadbeef bval=0x00000000   bit: 0xdeadbeef
[C] logic: aval=0x0000beef bval=0x00000000   bit: 0x0000beef
[SV] note from C: C reference active, K = 1.646760254
[SV] worst |model - C reference| = 3.597 LSB over 361 angles
\end{txtcode}

\begin{gotcha}{\kw{logic} and \kw{bit} are different C types}
The two \code{c\_show\_types} calls in the run above passed
\code{32'hDEAD\_BEEF} and then \code{32'hXXXX\_BEEF}. A \kw{logic} argument
arrives as an \code{svLogicVecVal}, a struct with an \code{aval} and a
\code{bval} word: \code{bval} bits mark the positions that are \code{X} or
\code{Z}. A \kw{bit} argument arrives as a plain \code{svBitVecVal}, one
word, with no way to represent \code{X}. If you declare the C parameter as
\code{svBitVecVal*} but the \SV{} side passes a \kw{logic}, you will read the
\code{aval} word and silently treat every \code{X} as whatever the simulator
happened to leave there.

Note also what \tool{Verilator} printed: \code{bval} is zero in both calls,
and the second call shows \code{aval=0x0000beef}. \tool{Verilator} is a
two-state simulator, so the \code{X} nibbles became zeros before the call was
made. On a four-state simulator the second line would show a non-zero
\code{bval}. Do not use \tool{Verilator} to test four-state DPI behaviour.
\end{gotcha}

\subsection{The same thing on \tool{QuestaSim}}

\index{QuestaSim}
\tool{QuestaSim} is not installed on the machine this chapter was written on,
so the following command lines are given as reference and are \emph{not}
verified here. They are the standard three steps: compile the \SV{}, build a
shared object, load it at elaboration.

\begin{shcode}
# Not executed for this book -- QuestaSim was unavailable.
vlib work
vlog -sv sv/cordic_pkg.sv sv/tb_dpi.sv
gcc -shared -fPIC -I$QUESTA_HOME/include \
    cpp/cordic_dpi.c -o cordic_dpi.so -lm
vsim -c -sv_lib cordic_dpi tb_dpi -do "run -all; quit"
\end{shcode}

Note that \code{-sv\_lib} takes the library name \emph{without} the
\file{.so} suffix; passing \code{cordic\_dpi.so} is the most common first
mistake. \code{vsim} also accepts \code{-dpicpppath} to select the compiler
it uses when it builds the shim itself.

\subsection{Embedding CPython in the shim}

\index{CPython!embedding}
Now the step that motivates this section: instead of reimplementing the model
in C, have the C shim call the Python one. CPython can be embedded in any
process. \Cref{lst:python:pydpi} is a working shim.

\begin{cppcode}[caption={A DPI shim with CPython embedded
(\file{py\_dpi.cpp}).},label={lst:python:pydpi}]
#define PY_SSIZE_T_CLEAN
#include <Python.h>
#include "svdpi.h"

extern "C" {

static PyObject* g_fn = nullptr;   /* cordic_model.cordic_rot */

void py_ref_open(const char* module, const char* func) {
  if (!Py_IsInitialized()) {
    PyConfig cfg;
    PyConfig_InitPythonConfig(&cfg);
    Py_InitializeFromConfig(&cfg);
    PyConfig_Clear(&cfg);
    PyRun_SimpleString("import sys; sys.path.insert(0, 'py')");
  }
  PyObject* m = PyImport_ImportModule(module);
  if (!m) { PyErr_Print(); return; }
  g_fn = PyObject_GetAttrString(m, func);
  Py_DECREF(m);
  if (!g_fn || !PyCallable_Check(g_fn)) { PyErr_Print(); g_fn = nullptr; }
}

void py_ref_call(int x0, int y0, int th, int* xo, int* yo) {
  *xo = 0; *yo = 0;
  if (!g_fn) return;
  PyGILState_STATE gil = PyGILState_Ensure();   /* one interpreter, one GIL */
  PyObject* r = PyObject_CallFunction(g_fn, "iii", x0, y0, th);
  if (!r) { PyErr_Print(); PyGILState_Release(gil); return; }
  *xo = (int)PyLong_AsLong(PyTuple_GetItem(r, 0));
  *yo = (int)PyLong_AsLong(PyTuple_GetItem(r, 1));
  Py_DECREF(r);
  PyGILState_Release(gil);
}

void py_ref_close(void) { Py_XDECREF(g_fn); g_fn = nullptr; Py_FinalizeEx(); }

}  /* extern "C" */
\end{cppcode}

Built against the system CPython and run under \tool{Verilator}:

\begin{shcode}
$ verilator --binary --timing --top-module tb_pydpi \
      -CFLAGS "$(python3.11-config --includes)" \
      -LDFLAGS "$(python3.11-config --embed --ldflags)" \
      sv/cordic_pkg.sv sv/tb_pydpi.sv cpp/py_dpi.cpp \
      -o Vtb_pydpi --Mdir obj_pydpi
$ ./obj_pydpi/Vtb_pydpi
\end{shcode}

\begin{txtcode}
[SV] theta=0 -> python says (16383, 1)
[SV] theta=2458 -> python says (15651, 4841)
[SV] theta=4915 -> python says (13521, 9251)
[SV] theta=7373 -> python says (10183, 12834)
[SV] theta=9830 -> python says (5936, 15270)
\end{txtcode}

It works. Now the costs, measured on the same machine, 50\,000 calls of each:

\begin{txtcode}
[bench] 50000 calls: embedded Python 0.818 s, C 0.018 s, ratio 46.7
\end{txtcode}

About \SI{16}{\micro\second} per Python call against \SI{0.36}{\micro\second}
for the C one -- between one and two orders of magnitude, and that is for a
model whose body is a fourteen-iteration integer loop. The gap is mostly
call overhead: argument marshalling, tuple allocation, the GIL. A model doing
real \code{numpy} work per call would look better in relative terms and worse
in absolute ones.

The other costs are not measurable but matter more:

\begin{itemize}
\item \textbf{Interpreter start-up} happens once, but it happens inside the
      simulator, and if it fails the simulator dies with a stack trace nobody
      can read.
\item \textbf{The GIL.} One interpreter per process. If your simulator is
      multi-threaded (\tool{Verilator} with \code{--threads}, or a
      commercial simulator running a multi-core partition) every DPI call
      into Python serialises.
\item \textbf{One process, one crash.} A segfault in a C extension module
      takes the simulation with it, mid-run, with no waveform flushed.
\item \textbf{Fragile builds.} The simulator, the Python interpreter and
      every extension module must agree on the C++ ABI and the C runtime.
      Mixing a conda Python with a vendor simulator's bundled libstdc++ is a
      well-known way to lose an afternoon.
\end{itemize}

\subsection{The pragmatic alternative}

If you already have the Python model and you want it in the loop, a
long-lived Python worker process talking to the simulator over a pipe or a
socket avoids every one of those problems except speed. The shim in C stays
tiny, the interpreter lives in its own process with its own crash domain, and
you can restart or debug the model independently. That is coupling style (d),
and \cref{sec:python:socket} shows it.

\begin{ruleofthumb}
Do not embed CPython in a simulator for a student project or a thesis. Use
offline vectors, generated source, or cocotb. Embed CPython only when you
have an existing Python model that must be in the loop, no cocotb-supported
simulator, and a build engineer.
\end{ruleofthumb}

\section{In-process, part 2: \tool{Verilator} and \code{pybind11}}
\label{sec:python:pybind}

\index{pybind11}
\tool{Verilator} compiles your design into a C++ class. A C++ class can be
wrapped for Python with \code{pybind11} in about forty lines. The result is a
testbench written entirely in Python, driving a compiled model directly, with
no VPI, no DPI and no simulator kernel in between.

\begin{cppcode}[caption={A \code{pybind11} wrapper around the Verilated model
(\file{pybind\_cordic.cpp}).},label={lst:python:pybind}]
#include <pybind11/pybind11.h>
#include "Vcordic_rot.h"
#include "verilated.h"

namespace py = pybind11;

class Cordic {
 public:
  Cordic() : ctx_(new VerilatedContext), dut_(new Vcordic_rot(ctx_)) {}
  ~Cordic() { dut_->final(); delete dut_; delete ctx_; }

  // One half cycle: set the clock, evaluate, advance simulation time.
  void half_tick(int level) {
    dut_->clk_i = level;
    dut_->eval();
    ctx_->timeInc(5000);   // 5 ns; timeInc counts timeprecision units
  }
  void tick() { half_tick(0); half_tick(1); }

  void set_rst(int v)   { dut_->rst_ni = v; }
  void set_valid(int v) { dut_->req_valid_i = v; }
  void set_ready(int v) { dut_->rsp_ready_i = v; }
  void set_req(int x, int y, int th) {
    dut_->req_x_i     = static_cast<uint16_t>(x);
    dut_->req_y_i     = static_cast<uint16_t>(y);
    dut_->req_theta_i = static_cast<uint16_t>(th);
  }
  void eval() { dut_->eval(); }
  bool req_ready() const { return dut_->req_ready_o; }
  bool rsp_valid() const { return dut_->rsp_valid_o; }
  int  rsp_x() const { return static_cast<int16_t>(dut_->rsp_x_o); }
  int  rsp_y() const { return static_cast<int16_t>(dut_->rsp_y_o); }
  uint64_t time() const { return ctx_->time(); }

 private:
  VerilatedContext* ctx_;
  Vcordic_rot* dut_;
};

PYBIND11_MODULE(vcordic, m) {
  m.doc() = "Verilated cordic_rot, driven from Python";
  py::class_<Cordic>(m, "Cordic")
      .def(py::init<>())
      .def("tick", &Cordic::tick).def("eval", &Cordic::eval)
      .def("set_rst", &Cordic::set_rst)
      .def("set_valid", &Cordic::set_valid)
      .def("set_ready", &Cordic::set_ready)
      .def("set_req", &Cordic::set_req)
      .def_property_readonly("req_ready", &Cordic::req_ready)
      .def_property_readonly("rsp_valid", &Cordic::rsp_valid)
      .def_property_readonly("rsp_x", &Cordic::rsp_x)
      .def_property_readonly("rsp_y", &Cordic::rsp_y)
      .def_property_readonly("time", &Cordic::time);
}
\end{cppcode}

\begin{tipbox}{Building the extension}
\tool{Verilator} must be told to compile position-independent code, because
the result goes into a shared object. Forgetting \code{-CFLAGS -fPIC} gives
a linker error about \code{R\_X86\_64\_TPOFF32} relocations that looks far
more alarming than it is.
\begin{shcode}
$ verilator --cc --top-module cordic_rot -CFLAGS -fPIC \
      sv/cordic_pkg.sv sv/cordic_rot.sv --Mdir obj_pb
$ make -C obj_pb -f Vcordic_rot.mk
$ g++ -O2 -std=c++17 -shared -fPIC \
      $(python3 -m pybind11 --includes) \
      -I obj_pb -I /usr/share/verilator/include \
      -I /usr/share/verilator/include/vltstd \
      cpp/pybind_cordic.cpp obj_pb/libVcordic_rot.a obj_pb/libverilated.a \
      -o vcordic.so
\end{shcode}
\end{tipbox}

The driver is then ordinary Python, with \code{numpy} available and the
golden model one import away.

\begin{pycode}[caption={Driving the Verilated model from Python
(\file{drive\_pybind.py}).},label={lst:python:pybinddrv}]
import numpy as np
import cordic_model as m
import vcordic

def run(dut, x0, y0, th):
    dut.set_req(x0, y0, th)
    dut.set_valid(1)
    dut.eval()
    while not dut.req_ready:
        dut.tick()
    dut.tick()                       # the accepting edge
    dut.set_valid(0)
    dut.eval()
    while not dut.rsp_valid:
        dut.tick()
    return dut.rsp_x, dut.rsp_y

dut = vcordic.Cordic()
dut.set_rst(0); dut.set_ready(1)
for _ in range(3):
    dut.tick()
dut.set_rst(1)

angles = np.linspace(-np.pi, np.pi, 721)
thetas = [m.float_to_q(a, m.QA, m.AW) for a in angles]
got = [run(dut, m.X0_UNIT, 0, th) for th in thetas]
exp = [m.cordic_rot(m.X0_UNIT, 0, th) for th in thetas]
print("bit-exact misses  :",
      sum(1 for g, e in zip(got, exp) if g != e))

gx = np.array([g[0] for g in got]) / float(1 << m.QF)
gy = np.array([g[1] for g in got]) / float(1 << m.QF)
err = np.maximum(np.abs(gx - np.cos(angles)), np.abs(gy - np.sin(angles)))
print(f"worst |err| vs sin/cos : {err.max() * (1 << m.QF):.2f} LSB")
\end{pycode}

\begin{txtcode}
$ python3 py/drive_pybind.py
transactions      : 721
bit-exact misses  : 0
worst |err| vs sin/cos : 3.82 LSB
simulated time    : 115380000 ps  (115.38 us)
wall time         : 8.0 ms (11.1 us per transaction)
\end{txtcode}

\code{ctx\_->time()} counts \code{timeprecision} units, and the design declares
\code{`timescale 1ns/1ps}, so the raw number is picoseconds. That is why
\code{half\_tick()} advances by 5000 and not by 5: see the discussion in
\cref{ch:verilator}.

\index{Verilator!limitations}
This is attractive and it is genuinely useful for algorithmic blocks. The
whole testbench is Python; \code{numpy} and \code{matplotlib} are right
there; there is no VPI, no DPI, no simulator licence.

The limits are equally real. You are restricted to \tool{Verilator}'s
synthesisable subset, so no \code{initial} blocks with delays, no classes, no
constrained random inside the design, and only partial four-state support.
You are writing the clocking yourself: every \code{eval()} and every
\code{timeInc()} is your responsibility, and a missing \code{eval()} after
changing an input is a bug that looks like a design bug. There is no
scheduler, so two "concurrent" drivers are two Python loops you have to
interleave by hand. And the wrapper is code you maintain: add a port to the
\DUT{} and you edit C++.

For a single-clock, valid/ready block this is fine. For anything with more
than two independent interfaces, use cocotb instead -- which gives you the
same Python, plus a scheduler.

\section{In-process, part 3: cocotb}
\label{sec:python:cocotb}

\index{cocotb}
cocotb (coroutine co-simulation testbench) is the option most students should
reach for. It is a Python library that drives a \DUT{} through the
simulator's own procedural interface -- VPI for Verilog, VHPI or FLI for
\VHDL{} -- and it replaces the HDL testbench entirely. There is no
\code{tb\_top.sv}. The top level of the simulation \emph{is} the \DUT{}, and
everything above it is Python.

\begin{figure}[htbp]
  \centering
  \begin{tikzpicture}[
      font=\small\sffamily,
      box/.style={draw,rounded corners=2pt,minimum height=9mm,
                  minimum width=34mm,align=center},
      ar/.style={-{Latex[length=2.2mm]},thick,draw=codenumber},
      ar2/.style={{Latex[length=2.2mm]}-{Latex[length=2.2mm]},thick,
                  draw=codenumber}]

    \node[box,draw=codekey,fill=notebg] (test)
      {\code{test\_cordic.py}\\[-1pt]{\scriptsize
       \code{@cocotb.test()} coroutines}};
    \node[box,draw=codekey,fill=notebg,below=7mm of test] (model)
      {\code{cordic\_model.py}\\[-1pt]{\scriptsize the golden model}};
    \node[box,draw=codekey,fill=notebg,right=16mm of test] (sched)
      {cocotb scheduler\\[-1pt]{\scriptsize triggers, tasks}};
    \node[box,draw=accent,fill=tipbg,below=7mm of sched] (gpi)
      {GPI\\[-1pt]{\scriptsize C++ abstraction layer}};
    \node[box,draw=accentalt,fill=svbg,below=7mm of gpi] (vpi)
      {VPI / VHPI / FLI};
    \node[box,draw=accentalt,fill=svbg,below=7mm of vpi] (sim)
      {simulator kernel\\[-1pt]{\scriptsize Questa, Verilator,
       Icarus, \dots}};
    \node[box,draw=accentalt,fill=svbg,below=7mm of sim] (dut)
      {\code{cordic\_rot}\\[-1pt]{\scriptsize the DUT -- the top level}};

    \draw[ar2] (test) -- (sched);
    \draw[ar]  (model) -- (test);
    \draw[ar2] (sched) -- (gpi);
    \draw[ar2] (gpi)   -- (vpi);
    \draw[ar2] (vpi)   -- (sim);
    \draw[ar2] (sim)   -- (dut);

    \node[right=2mm of sched.east,font=\scriptsize,align=left,
          text width=30mm]
      {\code{await}\\\code{RisingEdge(clk)}\\suspends a task and\\registers a
       callback};
    \node[right=2mm of vpi.east,font=\scriptsize,align=left,text width=30mm]
      {value read/write and callback registration, per IEEE 1800 clause 36};

    \begin{scope}[on background layer]
      \node[draw=codeframe,dashed,rounded corners=3pt,
            fit=(test)(model)(sched),inner sep=3mm,
            label={[font=\scriptsize\sffamily\bfseries,
                    text=codekey]above:Python process}] {};
    \end{scope}
  \end{tikzpicture}
  \caption{cocotb architecture. The Python test and the simulator run in one
  process; the scheduler suspends coroutines on simulator callbacks
  registered through the GPI, which is a thin C++ layer over VPI, VHPI or
  FLI. There is no HDL testbench.}
  \label{fig:python:cocotbarch}
\end{figure}

\begin{notebox}{Why there is no HDL testbench}
\index{VPI}\index{VHPI}
VPI (the Verilog Procedural Interface, IEEE 1800-2017 clause 36) and VHPI
(IEEE 1076-2008 clause 22) let an external C program walk the design
hierarchy, read and write signal values, and register callbacks on events
such as "value changed" or "end of time step". That is exactly what a driver
and a monitor need, so a testbench written against VPI does not need to be
instantiated inside the design at all. cocotb wraps VPI, VHPI and Mentor's
FLI behind one C++ abstraction called the GPI, and exposes it to Python. The
consequence for you is practical: the \code{TOPLEVEL} in a cocotb run is the
\DUT{} itself, and there is no HDL file to write.
\end{notebox}

\subsection{The core API}

\index{cocotb!triggers}
Everything in cocotb is a coroutine that awaits a trigger.

\begin{itemize}
\item \code{@cocotb.test()} marks a coroutine as a test. Every marked
      coroutine in the module named by \code{MODULE} is run, in order, and
      each gets the \DUT{} handle as its only argument.
\item \code{await Timer(10, unit="ns")} advances simulation time.
\item \code{await RisingEdge(dut.clk\_i)} suspends until the next rising
      edge. \code{FallingEdge}, \code{Edge} (either) and \code{ValueChange}
      are the siblings.
\item \code{await ClockCycles(dut.clk\_i, 5)} is five rising edges.
\item \code{cocotb.start\_soon(coro())} launches a background task and
      returns immediately -- this is how you get a driver and a monitor
      running concurrently. It is the cocotb equivalent of \kw{fork} \dots
      \kw{join\_none}.
\item \code{Clock(dut.clk\_i, 10, unit="ns").start()} is a coroutine that
      generates a clock; start it with \code{start\_soon}.
\item \code{await Combine(a, b)} waits for all of several tasks;
      \code{await First(a, b)} waits for the first of them --
      \kw{join} and \code{join\_any} respectively.
\item \code{await with\_timeout(task, 1000, "ns")} raises
      \code{SimTimeoutError} if the task has not finished in time. Use it on
      every driver: a hung handshake should fail the test, not hang the
      regression.
\item \code{dut.sig.value = 3} schedules a write; \code{dut.sig.value} reads.
      The read gives a \code{LogicArray} with \code{.integer} (unsigned),
      \code{.signed\_integer} (two's complement), \code{.binstr} (the string
      of \code{0}, \code{1}, \code{x} and \code{z}) and \code{.is\_resolvable}.
\end{itemize}

\subsection{A complete testbench for the CORDIC}

\Cref{lst:python:cocotbtb} is the whole thing: driver, monitor, scoreboard
and three tests, importing the golden model directly. Compare it with the
class-based \SV{} testbench of \refverificationbook: the structure is identical --
transaction, driver, monitor, scoreboard -- and the line count is about a
fifth.

\begin{pycode}[caption={A complete cocotb testbench for \code{cordic\_rot}
(\file{test\_cordic.py}).},label={lst:python:cocotbtb}]
import math, os, random, sys
from collections import deque
from pathlib import Path

import cocotb
from cocotb.clock import Clock
from cocotb.triggers import ClockCycles, RisingEdge, with_timeout

sys.path.insert(0, str(Path(__file__).resolve().parents[1] / "py"))
import cordic_model as m

TOL_LSB = 12


class Txn:
    """The transaction: stimulus in, observed response filled in later."""
    def __init__(self, x0, y0, theta):
        self.x0, self.y0, self.theta = x0, y0, theta
        self.x = self.y = None

    def __repr__(self):
        return f"Txn(x0={self.x0}, y0={self.y0}, theta={self.theta})"


async def reset(dut, cycles=5):
    dut.rst_ni.value = 0
    dut.req_valid_i.value = 0
    dut.rsp_ready_i.value = 1
    dut.req_x_i.value = dut.req_y_i.value = dut.req_theta_i.value = 0
    await ClockCycles(dut.clk_i, cycles)
    dut.rst_ni.value = 1
    await RisingEdge(dut.clk_i)


async def driver(dut, queue, stimulus):
    """Push transactions into the request channel, honouring req_ready_o."""
    for txn in stimulus:
        await RisingEdge(dut.clk_i)
        dut.req_x_i.value = txn.x0 & 0xFFFF
        dut.req_y_i.value = txn.y0 & 0xFFFF
        dut.req_theta_i.value = txn.theta & 0xFFFF
        dut.req_valid_i.value = 1
        while True:
            await RisingEdge(dut.clk_i)
            if dut.req_ready_o.value == 1:
                break
        dut.req_valid_i.value = 0
        queue.append(txn)


async def monitor(dut, queue, results, expected):
    """Sample the response channel on every accepted handshake."""
    while len(results) < expected:
        await RisingEdge(dut.clk_i)
        if dut.rsp_valid_o.value == 1 and dut.rsp_ready_i.value == 1:
            txn = queue.popleft()
            txn.x = dut.rsp_x_o.value.signed_integer
            txn.y = dut.rsp_y_o.value.signed_integer
            results.append(txn)


def scoreboard(results, log):
    """Bit-exact check against cordic_model, plus a float sanity check."""
    errors, worst = 0, 0.0
    for txn in results:
        ex, ey = m.cordic_rot(txn.x0, txn.y0, txn.theta)
        if (txn.x, txn.y) != (ex, ey):
            errors += 1
            if errors <= 5:
                log.error("%r: got (%d, %d) expected (%d, %d)",
                          txn, txn.x, txn.y, ex, ey)
        fx, fy = m.cordic_float(txn.x0, txn.y0, txn.theta / (1 << m.QA))
        worst = max(worst, abs(txn.x - fx), abs(txn.y - fy))
    return errors, worst


async def run(dut, stimulus):
    cocotb.start_soon(Clock(dut.clk_i, 10, unit="ns").start())
    await reset(dut)

    queue, results = deque(), []
    mon = cocotb.start_soon(monitor(dut, queue, results, len(stimulus)))
    drv = cocotb.start_soon(driver(dut, queue, stimulus))
    await with_timeout(drv, 200 * len(stimulus) + 1000, "ns")
    await with_timeout(mon, 200 * len(stimulus) + 1000, "ns")

    errors, worst = scoreboard(results, dut._log)
    dut._log.info("%d transactions, %d bit-exact mismatches, "
                  "worst float error %.2f LSB", len(results), errors, worst)
    assert errors == 0, f"{errors} bit-exact mismatches"
    assert worst <= TOL_LSB, f"worst float error {worst:.2f} > {TOL_LSB} LSB"


@cocotb.test()
async def test_directed(dut):
    """Angles a random generator will not hit."""
    angs = [0.0, math.pi/2, -math.pi/2, math.pi/4, -math.pi/4,
            1e-4, -1e-4, 3.14, -3.14]
    stim = [Txn(m.X0_UNIT, 0, m.float_to_q(a, m.QA, m.AW)) for a in angs]
    stim += [Txn(0, 0, 0), Txn(6000, -6000, 0)]
    await run(dut, stim)


@cocotb.test()
async def test_random(dut):
    """Constrained random stimulus, seeded from the environment."""
    seed = int(os.environ.get("SEED", "20250905"))
    rng = random.Random(seed)
    lim = m.float_to_q(math.pi, m.QA, m.AW)
    stim = [Txn(rng.randint(-6000, 6000), rng.randint(-6000, 6000),
                rng.randint(-lim, lim)) for _ in range(100)]
    dut._log.info("seed = %d", seed)
    await run(dut, stim)


# test_sweep is the same shape: 181 angles spread over [-pi, pi].
\end{pycode}

The build is a \file{Makefile} of a dozen lines. Everything after
\code{include} comes from cocotb.

\begin{makecode}
SIM           ?= icarus
TOPLEVEL_LANG ?= verilog

RTL := $(PWD)/../sv

ifeq ($(TOPLEVEL_LANG),verilog)
  VERILOG_SOURCES := $(RTL)/cordic_pkg.sv $(RTL)/cordic_rot.sv
  ifeq ($(SIM),verilator)
    EXTRA_ARGS += --timing -Wno-fatal
  endif
else
  VHDL_SOURCES := $(RTL)/cordic_pkg.vhd $(RTL)/cordic_rot.vhd
endif

TOPLEVEL := cordic_rot
MODULE   := test_cordic

include $(shell cocotb-config --makefiles)/Makefile.sim
\end{makecode}

Four variables carry all the configuration. \code{SIM} names the simulator;
\code{TOPLEVEL\_LANG} is \code{verilog} or \code{vhdl} and decides which of
\code{VERILOG\_SOURCES} and \code{VHDL\_SOURCES} is used; \code{TOPLEVEL} is
the name of the \DUT{} as the simulator sees it, which is where a
case-sensitivity mismatch between \VHDL{} and the makefile usually bites; and
\code{MODULE} is the Python module (not file) containing the tests. Two more
are worth knowing: \code{COMPILE\_ARGS} goes to the analysis step and
\code{EXTRA\_ARGS} to both analysis and elaboration, while \code{SIM\_ARGS}
goes only to the run. The language-revision switch is an \emph{analysis}
switch: \code{-2008} for \code{vcom} belongs in \code{COMPILE\_ARGS}, and a
switch that a tool needs at both analysis and elaboration belongs in
\code{EXTRA\_ARGS}. Put one of them in \code{SIM\_ARGS} by mistake and the run
either rejects it as an unknown option or, worse, looks for a library built to
a different revision and reports that it cannot find your entity --- a message
that points at the design and not at the command line.

\subsection{Running it}

With \tool{Verilator} as the back end:

\begin{shcode}
$ make -C cocotb_tb SIM=verilator
\end{shcode}

\begin{txtcode}
     0.00ns INFO  cocotb  Running on Verilator version 5.020 2024-01-01
     0.00ns INFO  cocotb  Running tests with cocotb v1.9.2
     0.00ns INFO  cocotb.regression  Found test test_cordic.test_directed
     0.00ns INFO  cocotb.regression  Found test test_cordic.test_random
     0.00ns INFO  cocotb.regression  Found test test_cordic.test_sweep
     0.00ns INFO  cocotb.regression  running test_directed (1/3)
  1820.00ns INFO  cocotb.cordic_rot  11 transactions, 0 bit-exact
                                     mismatches, worst float error 1.07 LSB
  1820.00ns INFO  cocotb.regression  test_directed passed
  1820.00ns INFO  cocotb.regression  running test_random (2/3)
  1820.00ns INFO  cocotb.cordic_rot  seed = 20250905
 17890.00ns INFO  cocotb.cordic_rot  100 transactions, 0 bit-exact
                                     mismatches, worst float error 2.06 LSB
 17890.00ns INFO  cocotb.regression  test_random passed
 46920.00ns INFO  cocotb.cordic_rot  181 transactions, 0 bit-exact
                                     mismatches, worst float error 3.60 LSB
 46920.00ns INFO  cocotb.regression  test_sweep passed
 ** TEST                        STATUS  SIM TIME (ns)  REAL TIME (s)  RATIO **
 ** test_cordic.test_directed    PASS        1820.00           0.02   79492 **
 ** test_cordic.test_random      PASS       16070.00           0.18   89996 **
 ** test_cordic.test_sweep       PASS       29030.00           0.32   90580 **
 ** TESTS=3 PASS=3 FAIL=0 SKIP=0            46920.00           0.54   87526 **
\end{txtcode}

Now the part that matters most to a \VHDL{} student. The \emph{same Python
file}, not one line changed, runs against the \VHDL{} implementation of the
same core. Only the makefile variables change:

\begin{shcode}
$ make -C cocotb_tb SIM=questa TOPLEVEL_LANG=vhdl COMPILE_ARGS=-2008
\end{shcode}

\begin{txtcode}
  1820.00ns INFO  cocotb.cordic_rot  11 transactions, 0 bit-exact
                                     mismatches, worst float error 1.07 LSB
  1820.00ns INFO  cocotb.regression  test_directed passed
 17890.00ns INFO  cocotb.cordic_rot  100 transactions, 0 bit-exact
                                     mismatches, worst float error 2.06 LSB
 17890.00ns INFO  cocotb.regression  test_random passed
 46920.00ns INFO  cocotb.cordic_rot  181 transactions, 0 bit-exact
                                     mismatches, worst float error 3.60 LSB
 46920.00ns INFO  cocotb.regression  test_sweep passed
 ** TESTS=3 PASS=3 FAIL=0 SKIP=0                                            **
\end{txtcode}

Same test names, same transaction counts, same bit-exact result, same worst
error. The wall-clock and the simulated-nanoseconds-per-second figures are the
one thing that does change, and they change with the simulator rather than
with the test, which is why they are left out above.

That is the argument in one line: you can use cocotb on your \VHDL{}
coursework today, and the testbench you write will still be valid when the
design is rewritten in \SV{}. cocotb officially supports
\tool{Questa}/\tool{ModelSim}, \tool{Xcelium}, \tool{VCS},
\tool{Riviera-PRO}, \tool{Verilator} and \tool{Icarus Verilog}, so whichever
tool your department has, the Python side is portable.

\begin{gotcha}{\code{.value} on a multi-bit signal is unsigned}
Reading a port gives you a bit pattern, and cocotb interprets it as an
unsigned integer by default -- regardless of whether the \SV{} declared it
\kw{logic} \kw{signed} or the \VHDL{} declared it \code{signed}. The extra
test below reads \code{rsp\_x\_o} after a rotation by \(\pi\), where the
answer should be \(-16384\):
\begin{txtcode}
binstr           = 1100000000000000
int(v)           = 49152
v.integer        = 49152
v.signed_integer = -16384
model says       = -16384
\end{txtcode}
Use \code{.signed\_integer} for any signed port; if you use \code{.integer}
or \code{int()}, every negative result will read as a large positive one and
your scoreboard will report a hundred per cent failure rate on a working
design. When writing, the reverse applies: \code{dut.req\_x\_i.value = -1234}
is accepted by recent cocotb, but masking explicitly
(\code{value \& 0xFFFF}) is unambiguous and works everywhere.
\end{gotcha}

\begin{gotcha}{cocotb and simulator versions are coupled}
cocotb talks to the simulator through VPI/VHPI, and both sides evolve. The
first attempt at the run above used cocotb 2.1.0 and failed immediately:
\begin{txtcode}
Makefile.verilator:29: *** cocotb requires Verilator 5.036 or later,
                           but using 5.020.  Stop.
\end{txtcode}
The run shown was produced with cocotb 1.9.2, which supports
\tool{Verilator} 5.020. Note also that cocotb 2.0 changed the API in ways
that matter for this chapter: \code{units=} became \code{unit=}, and the
value accessors were reorganised. Pin cocotb in your
\file{requirements.txt} together with the simulator version you tested
against, and write the pair into the CI job.
\end{gotcha}

\subsection{The ecosystem}

\index{pyuvm}\index{cocotb!cocotb-bus}\index{cocotb!cocotb-coverage}
cocotb is deliberately small. Four packages sit on top of it and you should
know they exist.

\code{cocotb-test} (and, in cocotb 2.x, the built-in \code{pytest} support)
replaces the \file{Makefile} with a \code{pytest} entry point. That gives you
\code{pytest -k sweep}, parameterised tests, fixtures, parallel runs with
\code{pytest-xdist} and a JUnit XML report that every CI system understands.
For a regression of more than a handful of tests it is worth the change.

\code{cocotb-bus} provides the \code{Driver}, \code{Monitor} and
\code{Scoreboard} base classes that used to live in cocotb itself, plus
ready-made verification IP for AXI4, AXI4-Lite, AXI-Stream, Avalon and
Wishbone. If your \DUT{} has an AXI port, do not write the driver.

\code{cocotb-coverage} adds constrained random generation and functional
coverage in Python: \code{@CoverPoint} and \code{@CoverCross} decorators that
collect bins as your tests run, and a constraint solver for randomisation.
It covers the same ground as \SV{}'s \kw{covergroup} and \kw{constraint}
blocks, in a form that will look familiar after this chapter.

\code{pyuvm} is UVM 1.2 reimplemented in Python: \code{uvm\_component},
\code{uvm\_sequence}, \code{uvm\_config\_db}, the factory, the phases, TLM
ports, all of it, running on cocotb. If you have read \refverificationbook, pyuvm
will be immediately legible, and it is an excellent way to learn UVM without
also fighting \SV{} class syntax and a licence server. It is not, today, what
industry signs off with.

\subsection{Honest limits}

\index{cocotb!performance}
\textbf{Speed.} Every \code{await RisingEdge(...)} is a callback from the
simulator into Python, and a testbench with a driver and a monitor both
waiting on the clock costs two Python round-trips per cycle. On this machine
the cocotb run above sustained about \SI{97000}{ns} of simulated time per
wall-clock second. The pure C++ \tool{Verilator} testbench of
\cref{sec:python:vectors}, driving the same design over the same handshake,
sustained about \SI{42}{\milli\second} of simulated time per second -- some
two to three orders of magnitude more. That ratio is workload dependent: a
testbench that awaits a transaction rather than a clock edge, on a design
with a hundred times more logic, will show a much smaller gap, because the
Python overhead is per callback and not per gate. Measure before you conclude
anything.

\textbf{No SVA.} You cannot write a concurrent assertion in Python. Assertions
that belong to the design should stay in the design, as in
\file{cordic\_rot.sv}, and they will still fire during a cocotb run under a
simulator that supports them. \tool{Verilator} supports immediate assertions
and a useful subset of concurrent ones; \tool{Icarus} does not parse
\kw{property} at all.

\textbf{Waveform debug is unchanged.} cocotb does not give you a Python
waveform viewer. You dump a VCD or FST from the simulator
(\code{EXTRA\_ARGS += --trace} for \tool{Verilator}; a \code{do} file with
\code{log -r /*} for \tool{QuestaSim}) and open it in the simulator's own wave
window, or in \tool{GTKWave} or \tool{Surfer} if it is a VCD or an FST ---
see \cref{sec:verilator:viewers}. The Python side is debugged with
\code{logging}, \code{pdb} or \code{pytest}.

\textbf{Industry.} Sign-off flows at the companies most of you will join are
\SV{} and UVM, driven by \tool{Questa}, \tool{VCS} or \tool{Xcelium}. cocotb
is widely used for block-level work, for IP developed by small teams, and
throughout the open-source silicon community, and it is growing. It is not a
reason to skip \refverificationbook.

\section{Out-of-process co-simulation}
\label{sec:python:socket}

\index{co-simulation!socket}
The last style is a small C shim inside the simulator that talks to a Python
process over a UNIX socket or a pipe. Use it when the model cannot be linked
in: a licensed library that only runs in its own process, a web service, a
physical model that takes seconds per call and must not block your build, or
simply a Python environment you cannot reconcile with the simulator's.

The shim is short. \Cref{lst:python:sock} is the whole client side; the
server is thirty lines of Python around \code{socket.AF\_UNIX}.

\begin{cppcode}[caption={A DPI shim over a UNIX socket
(\file{sock\_dpi.cpp}).},label={lst:python:sock}]
static int g_fd = -1;

int sock_open(const char* path) {
  struct sockaddr_un a;
  g_fd = socket(AF_UNIX, SOCK_STREAM, 0);
  if (g_fd < 0) return -1;
  std::memset(&a, 0, sizeof a);
  a.sun_family = AF_UNIX;
  std::snprintf(a.sun_path, sizeof a.sun_path, "%s", path);
  if (connect(g_fd, (struct sockaddr*)&a, sizeof a) < 0) return g_fd = -1;
  return 0;
}

int sock_ref(int x0, int y0, int th, int* xo, int* yo) {
  char buf[128];
  int n = std::snprintf(buf, sizeof buf, "%d %d %d\n", x0, y0, th);
  if (g_fd < 0 || write(g_fd, buf, n) != n) return -1;
  n = read(g_fd, buf, sizeof buf - 1);        /* blocking: deadlock risk */
  if (n <= 0) return -1;
  buf[n] = 0;
  return (std::sscanf(buf, "%d %d", xo, yo) == 2) ? 0 : -1;
}

void sock_close(void) { if (g_fd >= 0) close(g_fd); g_fd = -1; }
\end{cppcode}

\begin{txtcode}
$ python3 py/ref_server.py /tmp/cordic.sock &
[srv] listening on /tmp/cordic.sock
$ ./obj_sock/Vtb_sock
[SV] k=0 -> server says (16383, 1) rc=0
[SV] k=1 -> server says (15089, 6381) rc=0
[SV] k=2 -> server says (11413, 11754) rc=0
[SV] k=3 -> server says (5936, 15270) rc=0
[srv] served 4 requests
\end{txtcode}

The failure modes are the classic ones and they will all happen to you.

\textbf{Deadlock.} The \code{read} above is blocking. If the server crashes,
the simulation hangs forever, and in CI that means a job that runs until the
runner's timeout. Always set \code{SO\_RCVTIMEO} or use \code{poll} with a
deadline, and fail the test on timeout.

\textbf{Ordering.} A stream socket has no message framing. Two requests
issued before the first reply is read will be concatenated in the buffer and
your \code{sscanf} will read the wrong pair. Either keep the protocol
strictly request-response (as above), or put an explicit length prefix on
every message.

\textbf{Start-up race.} The simulator must not connect before the server has
bound the socket. In a \file{Makefile} that means a wait loop on the socket
file, not \code{sleep 1}.

\textbf{Cleanup.} A stale socket file from a crashed run makes the next run
fail with a confusing error. Unlink it on start-up, and put the socket path
under a per-run temporary directory.

\begin{ruleofthumb}
Socket co-simulation is a last resort. If you can write the model in Python
and use cocotb, do that instead; if you can precompute the answers, use
offline vectors.
\end{ruleofthumb}

\section{Choosing}
\label{sec:python:choosing}

\Cref{tab:python:choose} maps the cases you will actually meet onto a style.

\begin{table}[htbp]
\centering
\caption{Which coupling style to use.}
\label{tab:python:choose}
\small
\begin{tabularx}{\linewidth}{@{}lXX@{}}
\toprule
\textbf{Situation} & \textbf{Use} & \textbf{Why} \\
\midrule
Homework testbench, one simulator, a few hundred vectors
  & (a) offline vectors
  & Nothing to install, the file is the evidence, works everywhere. \\
Lab exercise in \VHDL{}
  & cocotb, with \code{TOPLEVEL\_LANG} set to \code{vhdl}
  & The HDL testbench disappears; the golden model is imported directly. \\
Constant table or ROM the tools will not compute
  & (b) generated source
  & Integer literals are accepted by every tool. Add a \code{gen-check}
    job. \\
Thesis accelerator, \tool{Verilator} in CI, algorithmic block
  & cocotb, or pybind11 if you need raw speed
  & Full \code{numpy} in the scoreboard; no licence. \\
Long random regression, millions of transactions
  & (a) offline vectors, or a C++ testbench
  & Python per-cycle callbacks would dominate the run time. \\
Industrial block, \tool{Questa}, UVM sign-off
  & (a) for vectors, DPI-C for a C reference model
  & UVM is the deliverable; Python stays outside the simulation. \\
Existing Python or licensed model that must be in the loop
  & (d) sockets
  & Crash isolation and environment isolation are worth the pain. \\
\bottomrule
\end{tabularx}
\end{table}

\begin{ruleofthumb}
Start with offline vectors. Move to cocotb when the \DUT{} needs stimulus
that depends on its own output. Move to DPI-C only when a simulator or a
methodology forces you to.
\end{ruleofthumb}

The four styles are not exclusive, and a mature flow usually runs three of
them at once. The CORDIC project in this book generates its constants at
build time, drives a fast offline-vector regression of a million transactions
in a C++ \tool{Verilator} testbench overnight, and runs a small cocotb suite
on every push for the interactive protocol cases -- backpressure, reset in
the middle of a rotation, an angle out of range. Each style is doing what it
is good at, and none of them duplicates the golden model, because all three
read \file{cordic\_model.py}.

That last point is the one to take away from the whole chapter. The question
is not "which coupling style is best" but "where does the single source of
truth live". Put it in one Python file, make everything else import it or be
generated from it, and the four styles become four ways of shipping the same
model to the same design.

\section{Putting it in CI}
\label{sec:python:ci}

\index{continuous integration}
A verification flow that only runs on your laptop is not a flow. The
\file{Makefile} below is the real one used for this chapter: it lints,
regenerates vectors when the model or the generator changes, runs the
simulation, and post-processes the results in Python. Every target has been
executed.

\begin{makecode}
PYTHON ?= python3
SEED   ?= 20250905
COUNT  ?= 200
RTL    := sv/cordic_pkg.sv sv/cordic_rot.sv
VEC    := vec/cordic_vec.txt vec/cordic_vec.hex vec/cordic_vec.n
GEN    := gen/cordic_tbl.h gen/cordic_tbl_pkg.sv

.PHONY: all lint gen gen-check sim check cocotb

all: lint sim check

lint:
	verilator --lint-only -Wall --top-module cordic_rot $(RTL)

$(VEC): py/gen_vectors.py py/cordic_model.py
	$(PYTHON) py/gen_vectors.py --seed $(SEED) --count $(COUNT) --outdir vec

$(GEN): py/gen_tables.py py/cordic_model.py
	$(PYTHON) py/gen_tables.py gen

gen-check: gen
	$(PYTHON) py/gen_tables.py .gencheck
	diff -r gen .gencheck || \
	  (echo "generated files are stale: run 'make gen'"; exit 1)

obj_file/Vtb_file: $(RTL) sv/tb_vec_file.sv
	verilator --binary --timing --top-module tb_vec_file \
	  $(RTL) sv/tb_vec_file.sv -o Vtb_file --Mdir obj_file

sim: obj_file/Vtb_file $(VEC)
	./obj_file/Vtb_file +DIR=vec

check: sim
	$(PYTHON) py/check_results.py vec/cordic_res.txt --tol-lsb 12

cocotb:
	$(MAKE) -C cocotb_tb SIM=verilator
	$(MAKE) -C cocotb_tb SIM=questa TOPLEVEL_LANG=vhdl COMPILE_ARGS=-2008
\end{makecode}

\begin{txtcode}
$ make gen-check all
python3 py/gen_tables.py .gencheck
diff -r gen .gencheck || (echo "generated files are stale..."; exit 1)
verilator --lint-only -Wall --top-module cordic_rot sv/cordic_pkg.sv \
    sv/cordic_rot.sv
python3 py/gen_vectors.py --seed 20250905 --count 200 --outdir vec
wrote 212 vectors to vec/cordic_vec.{txt,hex}
./obj_file/Vtb_file +DIR=vec
[tb] 217 lines, 212 vectors, 0 mismatches
[tb] PASS
python3 py/check_results.py vec/cordic_res.txt --tol-lsb 12
vectors        : 212
saturated      : 2 (excluded)
worst |error|  : 2.655 LSB
rms   error    : 0.791 LSB
tolerance      : 12.0 LSB
outside tol    : 0
\end{txtcode}

The corresponding GitHub Actions job is short because the \file{Makefile}
does the work. Note the pinned versions: this is the whole point of the
exercise.

\begin{txtcode}
name: verify
on: [push, pull_request]

jobs:
  cordic:
    runs-on: ubuntu-24.04
    steps:
      - uses: actions/checkout@v4

      - name: Install tools
        run: |
          sudo apt-get update && sudo apt-get install -y verilator
          python3 -m pip install "cocotb==1.9.2" numpy matplotlib jinja2

      - name: Generated files are up to date
        run: make gen-check

      - name: Lint, vectors, simulation, numerical check
        run: make lint all

      - name: cocotb regression
        run: make cocotb

      - name: Archive vectors and results
        if: always()
        uses: actions/upload-artifact@v4
        with:
          name: cordic-vectors
          path: |
            vec/
            cocotb_tb/results.xml
\end{txtcode}

\begin{tipbox}{Make the build fail, not the log}
Every step above fails the job on a mismatch, because
\code{check\_results.py} returns a non-zero exit code when a vector is
outside tolerance and \code{make} stops on the first failing recipe. A
testbench that prints \code{FAIL} and exits zero is worse than no testbench:
it trains everybody to read logs, and nobody reads logs. Finish every
checking script with an explicit \code{sys.exit(1)}, and every HDL testbench
with \code{\$fatal} on error.
\end{tipbox}

\begin{notebox}{Archive the vectors, not just the log}
The \code{upload-artifact} step matters more than it looks. When a nightly
regression fails on seed 8842 and nobody can reproduce it, the vector file
that produced the failure -- header and all -- is the difference between a
five-minute fix and a week of guessing. Keep it for every failing run.
\end{notebox}

\section{Checklist}
\label{sec:python:checklist}

\begin{itemize}
\item Write the reference model in the \RTL{} language if it fits on one
      screen and uses only arithmetic. Write it in Python as soon as it needs
      arrays, a transform, a file or a plot.
\item Default to \textbf{offline vectors}. They work with every simulator,
      they leave evidence on disk, and they cost nothing at run time.
\item A vector file is hex, one transaction per line, with a \code{//} header
      carrying the seed, every parameter and the model revision. Print that
      header from the testbench at time zero.
\item Never call \code{\$readmemh} without an explicit address range. It does
      not tell you how many entries it filled, and a short file passes
      silently -- especially under \tool{Verilator}, where the unfilled
      entries are zero rather than \code{'x}.
\item Check the return value of \code{\$sscanf}, and remember that
      \code{\$feof} is not a look-ahead: check \code{\$fgets} as well.
      \VHDL{}'s \code{endfile} is a look-ahead, which is why you forget.
\item Hex vector files carry no sign. Every reader must sign-extend
      explicitly: \code{\$signed} in \SV{}, \code{signed(...)} in \VHDL{},
      \code{static\_cast<int16\_t>} in C++, an explicit mask in Python.
\item Compare bit-exact models with \code{==}. Compare a floating point
      reference with a derived tolerance, exclude the cases the reference
      cannot model (saturation), and never raise the tolerance to make a test
      pass.
\item Every generated artefact needs a Make rule whose prerequisites include
      the generator and everything it imports, plus a CI job that regenerates
      and diffs.
\item Emit a \kw{package}, not an \code{`include}. One elaboration, one copy,
      one name.
\item Generate a constant table in Python when any tool in your flow rejects
      the elaboration-time version. \tool{Verilator} 5.020 accepts
      \code{\$atan} in a constant function; \tool{Icarus} 12.0 does not.
\item In DPI-C: \kw{pure} means no side effects and no callback; \kw{context}
      is required to call an exported function or to use \code{svGetScope}.
      \kw{logic} maps to \code{svLogicVecVal} (\code{aval}+\code{bval}) and
      \kw{bit} to \code{svBitVecVal}. Give every DPI symbol C linkage.
\item Do not embed CPython in a simulator unless you have no alternative. It
      costs one to two orders of magnitude per call, one GIL, one crash
      domain and a fragile build.
\item Use \textbf{cocotb} when the stimulus must react to the \DUT{}, and
      for \VHDL{} coursework in particular: the HDL testbench disappears and
      the Python one survives a rewrite of the design into \SV{}.
\item In cocotb, read signed ports with \code{.signed\_integer}, never
      \code{.integer} or \code{int()}. Wrap every driver in
      \code{with\_timeout}.
\item Pin cocotb and the simulator version together. cocotb 1.9.2 works with
      \tool{Verilator} 5.020; cocotb 2.1.0 requires 5.036 or later and
      renamed \code{units=} to \code{unit=}.
\item Sockets are the last resort. If you use them, frame the messages, set a
      receive timeout, and unlink the socket file on start-up.
\end{itemize}
